\scanpage{65}
% archive foliation: 64
(\unsure{Ds.} la cat. dérivée) puisque $C^\bullet$ est une résolution $\Sigma$-injective de $\Psi^0$.

\figure{\begin{tikzpicture}[x=1cm,y=1cm]\draw (0,0) .. controls (-.5,-.6) and (-1.3,-.2) .. (-.9,.1) .. controls (-.5,.25) and (.8,.12) .. (1.4,.24);\end{tikzpicture}}

Nous considérons maintenant les réalisations $\tfrac12$ simpliciales $\unsure{\mathcal P|P}_*\supset Q_*\supset \unsure{\mathcal P|P}_{0*}\supset K_*\ \unsure{\text{etc.}}$ de $\unsure{\mathcal P|P}\supset Q\supset \unsure{\mathcal P|P}_0\supset K$. Sur $\unsure{\mathcal P|P}_*$ il y a la structure $\unsure{\otimes_k|\otimes_*}$ ou $\otimes_*$
($L_*\otimes_*M_*$ défini composante par composante
\[
(L_*\otimes M_*)_{[n]}=L_{[n]}\otimes M_{[n]}
\]
)
avec contraintes d'ass, unité ($\Psi^0$), et de comm, \inserted{\unsure{de bilinéar.}} \deleted{Il est quasi évident} \inserted{On note que $K_*$ est un} \deleted{(\unsure{bidirectionnel}) $\unsure{\mathcal P|P}^*,Q^*$,} ``co-idéal'', \deleted{que $\unsure{\mathcal P|P}_0$ est un \unsure{idéal}} \inserted{(i.e. si l'un des facteurs} $L_*,M_*$ est dans $K_*$, \unsure{l'autre aussi}), $\unsure{\mathcal P|P}_0$ et $Q$ sont stables par $\otimes$, car

\margin{\deleted{$\unsure{K}$} \unsure{Il s'agit} de $\unsure{\Psi_*}$ \unsure{qui sont} \unsure{contractiles} (\unsure{dualité})}

\diagram{\begin{tikzcd}[column sep=small,row sep=small]
\Psi^0_* \arrow[d,equal] & \otimes & \Psi^0_* \arrow[d,equal] & \simeq & \Psi^0_* \arrow[d,equal] & (\Rightarrow\text{ stabilité de }\unsure{\unsure{\mathcal P|P}_0})\\
\mathbf1_* && \mathbf1_* && \mathbf1_*
\end{tikzcd}}

et
\[
\Psi^i_*\otimes\Psi^j_*\simeq\Psi^{i+j}_*+\text{élém de }K
\]
\[
(\Rightarrow\text{ stabilité de }\unsure{\unsure{\mathcal P|P}_{1*}=Q})
\]

On transporte cette structure multiplicative sur $\unsure{\mathcal P|P}$ et ses sous-catégories (notation $*$, pour éviter des confusions) et sur autres réalisations $\unsure{\mathcal P|P}_\bullet$ et $\unsure{\mathcal P|P}^*,\unsure{\mathcal P|P}^\bullet$ (ces deux dernières contravariantes). On note $*$ sur $\unsure{\mathcal P|P}_\bullet$ et $\unsure{\mathcal P|P}^\bullet$, et $\unsure{\otimes|\mathbin{\circledast}}$ sur $\unsure{\mathcal P|P}^*$, car l'équivalence $\vee$ de

